% =============================================================================
% KrishokChat — Innovation Due-Diligence Dossier
% Compile with: xelatex main.tex  (run twice for TOC/refs)
% Bengali font: Noto Sans Bengali (installed).  Main font: Segoe UI / Arial.
% Figures are OPTIONAL: missing files render as blank placeholder boxes.
% =============================================================================
\documentclass[11pt,a4paper]{article}

\usepackage[margin=18mm]{geometry}
\usepackage{microtype}
\usepackage{graphicx}
\usepackage{booktabs}
\usepackage{tabularx}
\usepackage{array}
\usepackage{multirow}
\usepackage{xcolor}
\usepackage{tikz}
\usepackage{pgfplots}
\pgfplotsset{compat=1.17}
\usetikzlibrary{shapes.geometric,arrows.meta,positioning,calc,backgrounds,fit}
\usepackage[most]{tcolorbox}
\usepackage{fontspec}
\usepackage{polyglossia}
\usepackage{enumitem}
\usepackage{fancyhdr}
\usepackage{titlesec}
\usepackage{caption}
\usepackage{longtable}
\usepackage{adjustbox}
\usepackage[hidelinks]{hyperref}
\usepackage{bookmark}

% ---------- Fonts ----------
\setdefaultlanguage{english}
\setmainfont{Arial}
\newfontfamily\bengalifont[Script=Bengali,Renderer=HarfBuzz]{Noto Sans Bengali}
\newcommand{\bn}[1]{{\bengalifont #1}}

% ---------- Palette ----------
\definecolor{kcgreen}{HTML}{1E5E3A}
\definecolor{kcearth}{HTML}{B87333}
\definecolor{kcblue}{HTML}{2F6FB0}
\definecolor{kcpurple}{HTML}{6B4FA0}
\definecolor{kcred}{HTML}{B33A3A}
\definecolor{kcamber}{HTML}{C9902A}
\definecolor{kcgrey}{HTML}{7A6F62}
\definecolor{kcbg}{HTML}{FAF8F5}
\definecolor{kcpanel}{HTML}{F4F1EA}
\definecolor{kcink}{HTML}{2C241E}
\definecolor{kcline}{HTML}{E5DFD5}

% ---------- Section styling ----------
\titleformat{\section}{\color{kcgreen}\bfseries\Large}{}{0pt}{}
\titleformat{\subsection}{\color{kcearth}\bfseries\large}{}{0pt}{}
\titlespacing*{\section}{0pt}{4pt}{6pt}
\setlength{\parindent}{0pt}
\setlength{\parskip}{5pt}

% ---------- Header / footer ----------
\pagestyle{fancy}
\fancyhf{}
\renewcommand{\headrulewidth}{0pt}
\renewcommand{\footrulewidth}{0.4pt}
\fancyfoot[L]{\footnotesize\color{kcgrey}KrishokChat\,|\,Bangladesh Innovation Fair 2026}
\fancyfoot[R]{\footnotesize\color{kcgrey}Page \thepage}

% ---------- Helpers ----------
% Evidence badges
\newtcbox{\badgeM}{on line,colback=kcgreen!12,colframe=kcgreen,boxrule=0.4pt,
  arc=2pt,left=3pt,right=3pt,top=1pt,bottom=1pt,fontupper=\scriptsize\bfseries\color{kcgreen}}
\newtcbox{\badgeS}{on line,colback=kcblue!10,colframe=kcblue,boxrule=0.4pt,
  arc=2pt,left=3pt,right=3pt,top=1pt,bottom=1pt,fontupper=\scriptsize\bfseries\color{kcblue}}
\newtcbox{\badgeP}{on line,colback=kcamber!15,colframe=kcamber,boxrule=0.4pt,
  arc=2pt,left=3pt,right=3pt,top=1pt,bottom=1pt,fontupper=\scriptsize\bfseries\color{kcearth}}
\newcommand{\MEASURED}{\badgeM{MEASURED}}
\newcommand{\SIMULATED}{\badgeS{SIMULATED}}
\newcommand{\ESTABLISHED}{\badgeM{ESTABLISHED}}
\newcommand{\TESTED}{\badgeS{TESTED}}
\newcommand{\PROJECTED}{\badgeP{PROJECTED}}

% One-line page takeaway under a title
\newcommand{\takeaway}[1]{{\color{kcgrey}\itshape\large #1}\par\vspace{2pt}%
  {\color{kcline}\rule{\linewidth}{0.8pt}}\par\vspace{6pt}}

% Source footer for a page
\newcommand{\srcfoot}[1]{\par\vspace{4pt}{\color{kcline}\rule{\linewidth}{0.4pt}}\par
  {\footnotesize\color{kcgrey}#1}\par}

% Pull quote
\newtcolorbox{pullquote}{colback=kcpanel,colframe=kcgreen,boxrule=0pt,
  leftrule=4pt,arc=0pt,left=10pt,right=10pt,top=8pt,bottom=8pt}

% Big number block
\newcommand{\bignum}[2]{\begin{minipage}[t]{\linewidth}\centering
  {\fontsize{34}{36}\selectfont\bfseries\color{kcgreen}#1}\\[2pt]
  {\footnotesize\color{kcink}#2}\end{minipage}}

% Figure placeholder that auto-swaps to real asset when it exists
% #1 = path, #2 = height, #3 = caption
\newcommand{\figslot}[3]{%
  \begin{center}
  \IfFileExists{#1}{\includegraphics[width=\linewidth,height=#2,keepaspectratio]{#1}}%
  {\setlength{\fboxsep}{0pt}%
   \fcolorbox{kcline}{kcpanel}{\parbox[c][#2][c]{\linewidth}{\centering
     {\color{kcgrey}\small\ttfamily [ figure placeholder ]}\\[2pt]
     {\footnotesize\color{kcgrey}\ttfamily\detokenize{#1}}}}}%
  \\[3pt]{\footnotesize\color{kcgrey}#3}
  \end{center}}

% Capability / info card
\newtcolorbox{card}[1]{colback=white,colframe=kcgreen,boxrule=0.6pt,arc=3pt,
  left=6pt,right=6pt,top=5pt,bottom=5pt,title={\bfseries #1},
  coltitle=white,colbacktitle=kcgreen,fonttitle=\small\bfseries}

% Diagram node styles
\tikzset{
   tbox/.style={rectangle,rounded corners=2pt,draw,thick,align=center,
    font=\footnotesize,minimum height=8mm,inner sep=4pt,text width=30mm},
  inbox/.style={tbox,fill=kcblue!12,draw=kcblue},
  evbox/.style={tbox,fill=kcgreen!12,draw=kcgreen},
  aibox/.style={tbox,fill=kcpurple!12,draw=kcpurple},
  redbox/.style={tbox,fill=kcred!12,draw=kcred},
  ambox/.style={tbox,fill=kcamber!18,draw=kcearth},
  greybox/.style={tbox,fill=black!6,draw=kcgrey},
  fl/.style={-{Latex[length=2mm]},thick,draw=kcgrey},
}

\begin{document}
\sloppy

% ============================ PAGE 1 � COVER ============================
\thispagestyle{empty}
\begin{tikzpicture}[remember picture,overlay]
  \fill[kcgreen] (current page.north west) rectangle ([yshift=-14mm]current page.north east);
  \fill[kcearth] ([yshift=-14mm]current page.north west) rectangle ([yshift=-16mm]current page.north east);
\end{tikzpicture}

\vspace{6mm}
\begin{center}
{\fontsize{46}{50}\selectfont\bfseries\color{kcgreen}KRISHOKCHAT}\\[4mm]
{\Large\bfseries\color{kcink}From Field Problem to Deployable Agricultural Intelligence}\\[3mm]
{\large\color{kcearth}Current Status, Research Evidence, System Architecture and Deployment Pathway}\\[8mm]
\end{center}

\figslot{figures/cover_field.png}{78mm}{A Bengali-first agricultural advisory platform for farmers and the institutions that support them.}

\vspace{6mm}
\begin{center}
{\large\bfseries\color{kcgreen}Research \textrightarrow{} Evidence \textrightarrow{} Engineering \textrightarrow{} Validation \textrightarrow{} Deployment}\\[10mm]
{\color{kcgrey}\large Bangladesh \textbullet{} 2026}
\end{center}

\newpage
% ==================== PAGE 2 � EXECUTIVE OVERVIEW ====================
\section*{\color{kcgreen}What KrishokChat is today}
\takeaway{The product is already built; the next stage is field validation and institutional deployment.}

\begin{minipage}[t]{0.50\linewidth}
KrishokChat is a Bengali-first agricultural intelligence platform built to help
farmers obtain useful agricultural guidance without requiring specialist
terminology, constant access to an expert, or uninterrupted connectivity.

\vspace{4pt}
The platform brings together conversational agricultural question answering,
evidence-grounded retrieval, crop and disease image analysis, safety screening,
response verification, provenance, auditability and constrained-delivery pathways
in one system. The engineering direction is deliberately broader than a chatbot:
KrishokChat is developed as a reusable advisory platform that can support farmers
directly and, through the same infrastructure, extension officers, NGOs,
agribusinesses and public agricultural services.
\end{minipage}\hfill
\begin{minipage}[t]{0.46\linewidth}
\vspace{-2pt}
\begin{card}{Bengali interaction}\footnotesize Understands formal Bengali, farmer colloquial speech and Romanized Banglish.\end{card}\vspace{3pt}
\begin{card}{Agricultural evidence}\footnotesize Grounded in 2,946 semantic knowledge units from official publications.\end{card}\vspace{3pt}
\begin{card}{Crop \& disease vision}\footnotesize Photo becomes structured context for the advisory case.\end{card}\vspace{3pt}
\begin{card}{Safety \& verification}\footnotesize Fail-closed safety gate plus typed relational verification.\end{card}\vspace{3pt}
\begin{card}{Low-connectivity support}\footnotesize Cache, offline pathways and deterministic SMS delivery.\end{card}\vspace{3pt}
\begin{card}{Production-oriented infrastructure}\footnotesize Modular backend, audit trail, deployment and test layers.\end{card}
\end{minipage}

\srcfoot{Sources: KrishokChat v2 benchmark \ESTABLISHED{}; system repository \cite{krishokchat_repo}.}


\newpage
% ==================== PAGE 3 — WHY THE PROJECT STARTED ====================
\section*{\color{kcgreen}The project started with a field problem, not a model}
\takeaway{We began by listening to farmers, not by choosing a model.}

Our starting point was practical. We went to farmers, listened to the questions
they were already trying to solve, and observed the conditions under which
agricultural decisions are actually made. The problems did not arrive one at a
time. Language, disease recognition, treatment decisions, irrigation and soil
questions, access to expertise and network availability were often part of the
same situation.

\vspace{3pt}
That changed the design objective. We were not trying to build another
agricultural chatbot. We were trying to build an agricultural advisory system
that could remain useful when the farmer's language was informal, the information
was incomplete, the decision was safety-sensitive, and the connection was
unreliable.

\vspace{4pt}
\figslot{figures/field_origin.png}{58mm}{Field interviews with smallholder farmers in Rajshahi and Natore informed the design.}

\vspace{4pt}
\begin{pullquote}
{\large\bfseries\color{kcgreen}``The farmer's real problem is not simply lack of information. It is getting the right information, in the right form, at the moment a decision has to be made.''}
\end{pullquote}

\newpage
% ==================== PAGE 4 — THE PROBLEM ====================
\section*{\color{kcgreen}Why ordinary digital advice is not enough}
\takeaway{A useful advisory system must survive informal language, incomplete input, safety-sensitivity and weak connectivity at once.}

\begin{center}
\begin{tikzpicture}[node distance=6mm]
  \node[inbox,text width=42mm] (f) {\bfseries FARMER};
  \node[greybox,below=8mm of f,text width=95mm,align=left,minimum height=22mm] (traits)
    {\footnotesize
     \textbullet{} speaks colloquially \quad
     \textbullet{} may use Banglish / dialect\\
     \textbullet{} may not know the disease name \quad
     \textbullet{} may only have a photo\\
     \textbullet{} may ask an incomplete question \quad
     \textbullet{} may have weak connectivity};
  \node[evbox,below=8mm of traits,text width=60mm] (sys) {\bfseries ADVISORY SYSTEM};
  \node[greybox,below=6mm of sys,text width=60mm,fill=kcgreen!8,draw=kcgreen]
    (out) {\bfseries MUST STILL PRODUCE A SAFE, USEFUL RESPONSE};
  \draw[fl] (f) -- (traits);
  \draw[fl] (traits) -- (sys);
  \draw[fl] (sys) -- (out);
\end{tikzpicture}
\end{center}

\subsection*{Language}
Official agricultural documents are usually written differently from how farmers
ask questions.

\subsection*{Knowledge}
A plausible agricultural answer can still be wrong when chemical identity,
formulation, dose, timing or crop applicability is misbound.

\subsection*{Delivery}
A correct answer is not useful if it cannot reach the farmer under the available
connectivity conditions.

\srcfoot{Treatment advice can become unsafe through wrong chemical identity, formulation, dose, interval, PHI or crop applicability; network and message constraints are distinct last-mile failure surfaces \cite{cea_manuscript}.}

\newpage
% ==================== PAGE 5 — FIRST RESEARCH FOUNDATION ====================
\section*{\color{kcgreen}Before building the advisor, we built the evidence foundation}
\takeaway{The benchmark showed that agricultural language modeling needs grounded evidence. The product therefore began with evidence, not generation.}

Our first major research contribution was a provenance-traceable Bengali
agricultural benchmark designed specifically for agricultural knowledge,
treatment advice, safety behavior and structured reasoning.

\vspace{3pt}
The released KrishokChat v2 benchmark contains 85,979 core instances across four
tracks: General Knowledge QA, Treatment QA, Safety Refusal and Re-query, and
Table QA. The resource is built from 284 government agricultural publications
across 13 institutions and six Bengali dialects, with every benchmark instance
linked back to its originating semantic knowledge unit and source publication.
It deliberately separates what is often collapsed into a single ``agricultural
QA'' score: ordinary knowledge, safety-critical treatment advice,
refusal/clarification behavior, and reasoning over structured agricultural tables.

\vspace{6pt}
\begin{center}
\begin{tabularx}{\linewidth}{*{3}{>{\centering\arraybackslash}X}}
\bignum{284}{government publications} & \bignum{13}{institutions} & \bignum{6}{Bengali dialects}\\[10mm]
\bignum{2{,}946}{semantic knowledge units} & \bignum{85{,}979}{core benchmark instances} & \bignum{1{,}000}{real-world farmer queries}\\
\end{tabularx}
\end{center}

\vspace{4pt}
\begin{center}\footnotesize
\begin{tabular}{lr}
\toprule
\textbf{Track} & \textbf{Instances}\\
\midrule
General Knowledge QA & 28,993\\
Treatment QA & 11,224\\
Safety Refusal / Re-query & 20,112\\
Table QA & 25,650\\
\bottomrule
\end{tabular}
\end{center}

\srcfoot{Source: KrishokChat v2 benchmark \ESTABLISHED{} \cite{krishokchat_v2}.}

\newpage
% ==================== PAGE 6 — WHAT v2 TAUGHT US ====================
\section*{\color{kcgreen}Better language behavior does not automatically produce safe advice}
\takeaway{Evidence supplied to a generator is not the same thing as evidence enforced by the system.}

\begin{minipage}[t]{0.48\linewidth}
\textbf{Closed-book treatment correctness}\par\vspace{3pt}
\begin{center}
\begin{tikzpicture}
\begin{axis}[
  width=\linewidth,height=52mm,
  xbar, xmin=0, xmax=55,
  bar width=5pt,
  enlarge y limits=0.14,
  xlabel={\footnotesize Treatment Correct (\%)},
  symbolic y coords={Qwen-2.5-7B,LLaMA-3.1-8B,GPT-OSS-120B,Gemma-4-26B,Gemini-2.5-FL},
  ytick=data,
  nodes near coords,
  every node near coord/.append style={font=\scriptsize},
  tick label style={font=\scriptsize},
  xmajorgrids,
]
\addplot[fill=kcearth!70,draw=kcearth] coordinates {
  (13.29,Qwen-2.5-7B) (12.43,LLaMA-3.1-8B) (32.92,GPT-OSS-120B) (38.73,Gemma-4-26B) (43.64,Gemini-2.5-FL)};
\end{axis}
\end{tikzpicture}
\end{center}
\end{minipage}\hfill
\begin{minipage}[t]{0.48\linewidth}
\textbf{Even with oracle evidence}\par\vspace{3pt}
Treatment hallucination remains \textbf{4.05--7.00\%} across the evaluated
systems. Providing the right passage does not, by itself, prevent an unsafe
treatment claim.

\vspace{8pt}
\begin{pullquote}
{\bfseries\color{kcgreen}Evidence helps. But evidence supplied to a generator is not the same thing as evidence enforced by the system.}
\end{pullquote}
\end{minipage}

\srcfoot{Source: KrishokChat v2 closed-book and oracle-evidence evaluations \MEASURED{} \cite{krishokchat_v2}.}

\newpage
% ==================== PAGE 7 — FARMER LANGUAGE ====================
\section*{\color{kcgreen}The farmer does not speak in benchmark language}
\takeaway{A system can ace a clean question and still fail when the farmer does not name the crop.}

The benchmark's controlled tracks were not enough. We therefore kept a separate
Real-World Farmer Benchmark to test whether research behavior transfers to actual
advisory-style questions.

\vspace{3pt}
The 1,000-query benchmark was collected independently of benchmark construction.
Three hundred queries came from structured field interviews with smallholder
farmers in Rajshahi and Natore; the remaining queries came from agricultural
Facebook communities and a farmer question-and-answer portal. The evaluation
split contains 350 queries.

\vspace{3pt}
These questions are shorter, more colloquial and often less complete than
benchmark-generated questions. That difference matters because a system can
perform well on a clean agricultural question and still fail when the farmer does
not name the crop, describes a symptom indirectly, or uses local terminology.

\vspace{4pt}
\begin{center}
\begin{tabularx}{0.9\linewidth}{*{3}{>{\centering\arraybackslash}X}}
\bignum{1{,}000}{farmer queries collected} & \bignum{300}{structured field interviews} & \bignum{350}{evaluation split}\\
\end{tabularx}
\end{center}

\srcfoot{Source: KrishokChat v2 Real-World Farmer Benchmark \ESTABLISHED{} \cite{krishokchat_v2}.}

\newpage
% ==================== PAGE 8 — RETRIEVAL PAPER ====================
\section*{\color{kcgreen}Where does retrieval fail?}
\takeaway{Retrieval behavior changes with the form of the question; evidence is not equally reachable from every farmer query.}

Our second research study treated retrieval failure as its own problem rather
than hiding it behind an end-to-end chatbot score. Across 1,000 evaluation
queries and 2,882 knowledge nodes extracted from 284 official Bangladeshi
agricultural publications, we compared five retrieval architectures and six
embedding models under different language conditions.

\vspace{3pt}
The result was not that one retriever always wins. The stronger finding was that
retrieval behavior changes with the form of the question. Colloquial farmer
language, formal Bengali, and other language conditions create different failure
patterns, so an agricultural advisory system should not assume that text
retrieval is equally reliable for every farmer query. Hybrid RRF reaches overall
R@10 of 0.539, while dense retrieval falls sharply on colloquial farmer queries;
BM25 is strong for native Bengali and degrades heavily on English-to-Bengali
retrieval.

\vspace{4pt}
\begin{center}
\begin{tabularx}{0.9\linewidth}{*{4}{>{\centering\arraybackslash}X}}
\bignum{5}{retrieval architectures} & \bignum{6}{embedding models} & \bignum{2{,}882}{knowledge nodes} & \bignum{0.539}{Hybrid RRF R@10}\\
\end{tabularx}
\end{center}

\begin{pullquote}
{\bfseries\color{kcgreen}The retrieval study showed that evidence is not equally reachable from every farmer query. The system therefore needed more than one route to the same agricultural knowledge.}
\end{pullquote}

\srcfoot{Source: ``Where Does Retrieval Fail?'' \ESTABLISHED{} \cite{retrieval_study}.}

\newpage
% ==================== PAGE 9 — TURNING POINT ====================
\section*{\color{kcgreen}The research changed our architecture}
\takeaway{The architecture is the engineering response to the failure modes the research identified.}

\begin{minipage}[t]{0.44\linewidth}
\centering\textbf{What we initially had}\par\vspace{4pt}
\begin{tikzpicture}[node distance=5mm]
  \node[inbox,text width=32mm] (a) {Question};
  \node[greybox,below=of a,text width=32mm] (b) {Retrieve};
  \node[aibox,below=of b,text width=32mm] (c) {Generate};
  \node[evbox,below=of c,text width=32mm] (d) {Answer};
  \draw[fl](a)--(b);\draw[fl](b)--(c);\draw[fl](c)--(d);
\end{tikzpicture}
\end{minipage}\hfill
\begin{minipage}[t]{0.52\linewidth}
\centering\textbf{What the evidence demanded}\par\vspace{4pt}
\begin{tikzpicture}[node distance=3.2mm]
  \node[inbox,text width=44mm] (a) {Question / Image};
  \node[greybox,below=of a,text width=44mm] (b) {Understand context};
  \node[ambox,below=of b,text width=44mm] (c) {Assess safety \& uncertainty};
  \node[greybox,below=of c,text width=44mm] (d) {Route to the right capability};
  \node[evbox,below=of d,text width=44mm] (e) {Use structured evidence where possible};
  \node[aibox,below=of e,text width=44mm] (f) {Generate only where necessary};
  \node[evbox,below=of f,text width=44mm] (g) {Verify before delivery};
  \node[redbox,below=of g,text width=44mm] (h) {Answer / Clarify / Refuse / Escalate};
  \draw[fl](a)--(b);\draw[fl](b)--(c);\draw[fl](c)--(d);\draw[fl](d)--(e);
  \draw[fl](e)--(f);\draw[fl](f)--(g);\draw[fl](g)--(h);
\end{tikzpicture}
\end{minipage}

\vspace{4pt}
\begin{pullquote}
{\bfseries\color{kcgreen}The system architecture is not a collection of features added around a chatbot. It is the engineering response to the failure modes identified by the research program.}
\end{pullquote}

\newpage
% ==================== PAGE 10 — FULL ARCHITECTURE ====================
\section*{\color{kcgreen}One platform, one advisory case, multiple evidence paths}
\takeaway{The farmer meets one system; behind it, several evidence paths serve the same decision.}

\begin{center}
\begin{tikzpicture}[node distance=6mm and 10mm]
  \node[inbox] (farmer) {\bfseries FARMER};
  \node[inbox,below left=8mm and 6mm of farmer,text width=22mm] (text) {Text};
  \node[inbox,below=8mm of farmer,text width=22mm] (img) {Image};
  \node[inbox,below right=8mm and 6mm of farmer,text width=22mm] (voice) {Voice*};
  \draw[fl](farmer)--(text);\draw[fl](farmer)--(img);\draw[fl](farmer)--(voice);

  \node[ambox,below=20mm of farmer,text width=42mm] (ctx) {Context \& Safety};
  \draw[fl](text)--(ctx);\draw[fl](img)--(ctx);\draw[fl](voice)--(ctx);
  \node[greybox,below=of ctx,text width=42mm] (case) {Case State};
  \draw[fl](ctx)--(case);
  \node[greybox,below=of case,text width=42mm] (router) {Capability Router};
  \draw[fl](case)--(router);

  \node[evbox,below left=9mm and 4mm of router,text width=26mm] (facts) {Structured Facts};
  \node[evbox,below=9mm of router,text width=26mm] (rag) {RAG};
  \node[aibox,below right=9mm and 4mm of router,text width=26mm] (vision) {Vision Models};
  \draw[fl](router)--(facts);\draw[fl](router)--(rag);\draw[fl](router)--(vision);

  \node[evbox,below=26mm of router,text width=52mm] (evi) {Agricultural Evidence};
  \draw[fl](facts)--(evi);\draw[fl](rag)--(evi);\draw[fl](vision)--(evi);
  \node[greybox,below=of evi,text width=52mm] (resp) {Response Engine};
  \draw[fl](evi)--(resp);
  \node[greybox,below left=8mm and 2mm of resp,text width=30mm] (det) {Deterministic};
  \node[aibox,below right=8mm and 2mm of resp,text width=30mm] (llm) {LLM Realization};
  \draw[fl](resp)--(det);\draw[fl](resp)--(llm);
  \node[evbox,below=24mm of resp,text width=40mm] (verif) {Verification};
  \draw[fl](det)--(verif);\draw[fl](llm)--(verif);

  \node[evbox,below left=8mm and 6mm of verif,text width=22mm] (ans) {Answer};
  \node[ambox,below=8mm of verif,text width=22mm] (clar) {Clarify};
  \node[redbox,below right=8mm and 6mm of verif,text width=22mm] (esc) {Escalate};
  \draw[fl](verif)--(ans);\draw[fl](verif)--(clar);\draw[fl](verif)--(esc);
  \node[greybox,below=8mm of clar,text width=44mm] (deliver) {Web / Local / Cache / SMS};
  \draw[fl](clar)--(deliver);
\end{tikzpicture}
\end{center}

\srcfoot{\footnotesize *Voice shown as a planned input path. Colour key: {\color{kcblue}blue}=input/context, {\color{kcgreen}green}=verified evidence, {\color{kcpurple}purple}=generation, {\color{kcred}red}=safety stop, {\color{kcearth}amber}=uncertainty/clarification, {\color{kcgrey}grey}=infrastructure.}

\newpage
% ==================== PAGE 11 — FARMER JOURNEY ====================
\section*{\color{kcgreen}The system is designed around a real advisory interaction}
\takeaway{A farmer interaction is a case, not a single prompt.}

\begin{center}
\begin{tabularx}{\linewidth}{*{4}{>{\centering\arraybackslash}X}}
\figslot{screenshots/bengali_chat.png}{42mm}{} &
\figslot{screenshots/vision.png}{42mm}{} &
\figslot{screenshots/evidence.png}{42mm}{} &
\figslot{screenshots/agent_trace.png}{42mm}{}\\
{\bfseries\color{kcgreen}Ask} & {\bfseries\color{kcgreen}Understand} & {\bfseries\color{kcgreen}Verify} & {\bfseries\color{kcgreen}Act}\\
{\footnotesize Ask in Bengali} & {\footnotesize Add a crop image} & {\footnotesize Inspect evidence} & {\footnotesize Advice or clarification}\\
\end{tabularx}
\end{center}

\vspace{4pt}
KrishokChat treats a farmer interaction as a case rather than a single prompt.
The system can carry the farmer's question, available crop context, image
evidence, retrieved agricultural sources, verification status and subsequent
follow-up within the same advisory flow. This allows different capabilities to
contribute to the same decision instead of forcing the farmer to understand which
technical tool should be used first.

\srcfoot{Source: KrishokChat live system \cite{krishokchat_repo}.}

\newpage
% ==================== PAGE 12 — LANGUAGE LAYER ====================
\section*{\color{kcgreen}Built for the language farmers actually use}
\takeaway{The system is designed and evaluated for multiple Bengali registers and dialect conditions.}

\begin{center}
\begin{tcolorbox}[colback=kcpanel,colframe=kcline,width=0.9\linewidth,left=8pt,right=8pt]
\footnotesize
\textbf{Formal Bengali}\quad \bn{???? ??? ?????? ????? ??????????? ???}\\[3pt]
\textbf{Farmer-style Bengali}\quad \bn{???? ?????? ???? ??? ????? ?? ?????}\\[3pt]
\textbf{Banglish}\quad \texttt{alu te pata kalo hoye jacche ki korbo?}
\end{tcolorbox}
\end{center}

\vspace{3pt}
\begin{center}
\begin{tikzpicture}[node distance=5mm]
  \node[inbox,text width=40mm] (a) {Farmer expression};
  \node[ambox,right=of a,text width=44mm] (b) {Language normalization / interpretation};
  \node[greybox,right=of b,text width=32mm] (c) {Agricultural entity};
  \node[evbox,below=6mm of c,text width=32mm] (d) {Evidence retrieval};
  \draw[fl](a)--(b);\draw[fl](b)--(c);\draw[fl](c)--(d);
\end{tikzpicture}
\end{center}

\vspace{3pt}
Detection-gated coverage rises from \textbf{42.9\% to 89.8\%} for regional
dialects and \textbf{41.6\% to 90.0\%} for Romanized Banglish in the specified
evaluation. We do not claim that Bengali variation is solved; we claim that the
system is designed and evaluated for multiple Bengali registers and dialect
conditions, with safe fallback behavior when language becomes difficult.

\srcfoot{Source: KrishokChat dialect evaluation (E21) \MEASURED{}.}

\newpage
% ==================== PAGE 13 — MULTIMODAL DIAGNOSIS ====================
\section*{\color{kcgreen}A crop image becomes evidence for the same advisory case}
\takeaway{A photograph is not treated as an isolated classification task; it is structured context.}

\begin{center}
\begin{tikzpicture}[node distance=4.5mm]
  \node[inbox,text width=26mm] (p) {Photo};
  \node[aibox,below=of p,text width=40mm] (c) {Crop classification};
  \node[aibox,below=of c,text width=40mm] (d) {Crop-specific disease model};
  \node[ambox,below=of d,text width=40mm] (conf) {Confidence};
  \node[evbox,below=of conf,text width=40mm] (k) {Agricultural knowledge};
  \node[evbox,below=of k,text width=40mm] (adv) {Advisory};
  \draw[fl](p)--(c);\draw[fl](c)--(d);\draw[fl](d)--(conf);\draw[fl](conf)--(k);\draw[fl](k)--(adv);
\end{tikzpicture}
\end{center}

\vspace{3pt}
The vision workflow is deliberately integrated with the advisory layer. A
photograph is not treated as an isolated classification task; the predicted crop
and disease become structured context for downstream agricultural retrieval and
advice. This separation also allows uncertainty to propagate: a low-confidence or
contradictory visual prediction does not have to become a confident agricultural
recommendation. It can instead trigger clarification or another evidence path.

\srcfoot{Crop classification followed by per-crop disease routing; the vision output is treated as routing metadata, not a new detector \cite{krishokchat_repo,cea_manuscript}.}

\newpage
% ==================== PAGE 14 — CROSS-MODAL CONFLICT ====================
\section*{\color{kcgreen}When evidence disagrees, KrishokChat asks before it acts}
\takeaway{Conflicting evidence is itself a system state.}

\begin{center}
\begin{tcolorbox}[colback=kcamber!12,colframe=kcearth,width=0.8\linewidth,arc=3pt]
\centering\footnotesize
\textbf{TEXT:} Tomato \qquad\qquad \textbf{IMAGE:} Potato (0.91)\\[6pt]
{\large\bfseries\color{kcred}$\triangle$\ CONFLICT}\\[6pt]
``Your message and image suggest different crops. Please confirm before I continue.''\\[6pt]
{\footnotesize [\,Confirm text\,]\quad[\,Use image\,]\quad[\,Upload another image\,]}
\end{tcolorbox}
\end{center}

\vspace{3pt}
\begin{pullquote}
{\bfseries\color{kcgreen}A multimodal system should not silently choose whichever signal is more convenient. Conflicting evidence is itself a system state.}
\end{pullquote}

\vspace{3pt}
In the cross-modal-conflict evaluation, the system reached \textbf{100\%}
clarification triggering and \textbf{0\%} unsafe acceptance across the 100 tested
cases: a product feature, a safety feature, and a demonstrable system property at
once.

\srcfoot{Source: cross-modal conflict evaluation (E31) \MEASURED{}.}

\newpage
% ==================== PAGE 15 — SAFETY BOUNDARY ====================
\section*{\color{kcgreen}Safety is a system state, not a disclaimer}
\takeaway{An unsafe request has no normal generation path.}

\begin{center}
\begin{tikzpicture}[node distance=7mm and 12mm]
  \node[inbox,text width=40mm] (u) {User Input};
  \node[redbox,below=of u,text width=40mm] (gate) {Safety Gate};
  \draw[fl](u)--(gate);
  \node[evbox,below left=10mm and 2mm of gate,text width=26mm] (safe) {Safe};
  \node[ambox,below=10mm of gate,text width=26mm] (unclear) {Unclear};
  \node[redbox,below right=10mm and 2mm of gate,text width=26mm] (unsafe) {Unsafe};
  \draw[fl](gate)--(safe);\draw[fl](gate)--(unclear);\draw[fl](gate)--(unsafe);
  \node[greybox,below=6mm of safe,text width=26mm] (route) {Route};
  \node[greybox,below=6mm of unclear,text width=26mm] (clar) {Clarify};
  \node[redbox,below=6mm of unsafe,text width=26mm] (block) {Block};
  \draw[fl](safe)--(route);\draw[fl](unclear)--(clar);\draw[fl](unsafe)--(block);
\end{tikzpicture}
\end{center}

\begin{center}\textbf{\color{kcred}Unsafe \textrightarrow{} no normal generation path}\end{center}

\vspace{3pt}
KrishokChat treats safety-sensitive agricultural requests differently from
ordinary information requests. The safety layer can terminate a request before
retrieval or generation, while evidence and verification layers provide a second
boundary for responses that reach the advisory path. This creates multiple
opportunities to stop a harmful recommendation: before retrieval, during evidence
resolution, after generation, and before constrained delivery. Every refusal
surfaces the Bangladesh Krishi Call Center, \textbf{16123}.

\srcfoot{Six safety categories with terminal handling before retrieval and a local audit trail \cite{krishokchat_repo}.}

\newpage
% ==================== PAGE 16 — TREATMENT CONTRACT ====================
\section*{\color{kcgreen}An agricultural treatment is not a sentence --- it is a set of linked facts}
\takeaway{Correctness depends on relationships among fields, not the presence of individual words.}

\begin{minipage}[t]{0.42\linewidth}
\begin{tcolorbox}[colback=white,colframe=kcgreen,title={\bfseries Agricultural Treatment Record},
  coltitle=white,colbacktitle=kcgreen,fonttitle=\bfseries]
\footnotesize
Crop \\ Problem \\ Growth stage \\ Active ingredient \\ Formulation \\ Dose min \\
Dose max \\ Unit \\ Volume \\ Application interval \\ PHI \\ Regulatory status
\end{tcolorbox}
\end{minipage}\hfill
\begin{minipage}[t]{0.54\linewidth}
For safety-critical treatment advice, correctness depends on relationships among
fields rather than the presence of individual words.

\vspace{4pt}
A recommendation can contain the right chemical name and still be unsafe if the
dosage belongs to another crop, the formulation is different, the pre-harvest
interval is unsupported, or the recommendation comes from an obsolete or
conflicting source. KrishokChat therefore represents critical treatment
information as a typed relational record and requires the necessary fields to
remain jointly supportable before certification.
\end{minipage}

\srcfoot{Reflects the CEA treatment contract and its typed relational verification design \cite{cea_manuscript}.}

\newpage
% ==================== PAGE 17 — COUNTERFACTUAL BINDING ====================
\section*{\color{kcgreen}The verifier checks whether the facts belong together}
\takeaway{A word-matching system can certify a corrupted answer; a typed verifier does not.}

\begin{center}
\begin{tikzpicture}
\begin{axis}[
  width=0.9\linewidth,height=54mm,
  ybar, bar width=16pt,
  ymin=0,ymax=105,
  ylabel={\footnotesize Certified (\%)},
  symbolic x coords={Vanilla RAG,Lexical,LLM judge,Typed verifier},
  xtick=data,
  x tick label style={font=\scriptsize,align=center},
  nodes near coords,
  every node near coord/.append style={font=\scriptsize},
  legend style={font=\scriptsize,at={(0.5,-0.28)},anchor=north,legend columns=2},
  ymajorgrids,
]
\addplot[fill=kcgreen!60,draw=kcgreen] coordinates {(Vanilla RAG,100)(Lexical,100)(LLM judge,95.7)(Typed verifier,100)};
\addplot[fill=kcred!60,draw=kcred] coordinates {(Vanilla RAG,72.65)(Lexical,59.55)(LLM judge,38.65)(Typed verifier,0)};
\legend{Correct evidence certified,Counterfactual evidence certified}
\end{axis}
\end{tikzpicture}
\end{center}

\begin{center}{\Large\bfseries\color{kcgreen}CBC = 1.000}\end{center}

We tested the verification boundary by deliberately changing evidence that should
change a recommendation --- scaling the dose, shortening the pre-harvest interval,
or substituting incompatible evidence. A system that merely searches for matching
words can still certify the corrupted answer. In the 2,000-pair counterfactual
benchmark, the typed relational verifier certified \textbf{0.0\%} of corrupted
cases, while Vanilla RAG certified \textbf{72.65\%}.

\srcfoot{Source: 2,000-pair counterfactual benchmark \MEASURED{} \cite{cea_manuscript}.}

\newpage
% ==================== PAGE 18 — SLOT ABLATION ====================
\section*{\color{kcgreen}The safety schema changes behavior}
\takeaway{Removing different fields creates measurably different safety failures.}

\begin{center}
\begin{tikzpicture}
\begin{axis}[
  width=0.92\linewidth,height=72mm,
  xbar, bar width=5pt, xmin=0, xmax=35,
  xlabel={\footnotesize Dangerous-acceptance increase (percentage points)},
  symbolic y coords={Interval,PHI,Growth stage,Unit,Formulation,Volume denom.,Target pathogen,Host crop,Active ingredient,Regulatory polarity,Dosage bounds},
  ytick=data,
  nodes near coords,
  every node near coord/.append style={font=\scriptsize},
  tick label style={font=\scriptsize},
  xmajorgrids,
]
\addplot[fill=kcearth!70,draw=kcearth] coordinates {
  (1.9,Interval)(2.7,PHI)(4.9,Growth stage)(5.4,Unit)(5.8,Formulation)
  (7.1,Volume denom.)(8.2,Target pathogen)(11.4,Host crop)(14.2,Active ingredient)
  (17.8,Regulatory polarity)(31.6,Dosage bounds)};
\end{axis}
\end{tikzpicture}
\end{center}

Removing all typed constraints produces an \textbf{80.0\%} hazard rate; the full
tested configuration had \textbf{zero} observed dangerous acceptance across the
10,000 cases.

\begin{pullquote}
{\bfseries\color{kcgreen}The schema is not administrative metadata. Removing different fields creates measurably different safety failures, with dosage, regulatory status and active ingredient producing the largest individual effects.}
\end{pullquote}

\srcfoot{Source: 10,000-case slot-ablation study \MEASURED{} \cite{cea_manuscript}.}

\newpage
% ==================== PAGE 19 — RETRIEVAL DEGRADATION ====================
\section*{\color{kcgreen}When retrieval gets worse, the system refuses --- it does not invent}
\takeaway{Retrieval failure should reduce confidence, not increase the model's freedom to improvise.}

\begin{center}
\begin{tikzpicture}
\begin{axis}[
  width=0.9\linewidth,height=56mm,
  ybar, bar width=14pt, ymin=0, ymax=40,
  ylabel={\footnotesize Hazard (\%)},
  symbolic x coords={Gold,Noisy,Contradictory,Omitted},
  xtick=data, tick label style={font=\scriptsize},
  nodes near coords, every node near coord/.append style={font=\scriptsize},
  legend style={font=\scriptsize,at={(0.5,-0.22)},anchor=north,legend columns=2},
  ymajorgrids,
]
\addplot[fill=kcred!60,draw=kcred] coordinates {(Gold,10.0)(Noisy,19.2)(Contradictory,33.2)(Omitted,23.2)};
\addplot[fill=kcgreen!60,draw=kcgreen] coordinates {(Gold,0)(Noisy,0)(Contradictory,0)(Omitted,0)};
\legend{Vanilla RAG hazard,KrishokChat hazard}
\end{axis}
\end{tikzpicture}
\end{center}

This is one of the most important design decisions in the system. Retrieval
failure should reduce confidence in the answer, not increase the model's freedom
to improvise. We evaluate degradation as a first-class operating condition. As
authoritative evidence disappears or is replaced by contradictory material, the
system loses coverage through abstention rather than preserving a superficially
high answer rate by guessing.

\srcfoot{Source: retrieval-degradation evaluation \MEASURED{} \cite{cea_manuscript}.}

\newpage
% ==================== PAGE 20 — GENERATOR INDEPENDENCE ====================
\section*{\color{kcgreen}The safety layer is intended to outlive the model}
\takeaway{The system needs a model-independent boundary, not one perfect model.}

\begin{center}
\begin{tikzpicture}
\begin{axis}[
  width=0.9\linewidth,height=54mm,
  ybar, bar width=16pt, ymin=0, ymax=16,
  ylabel={\footnotesize Hazard (\%)},
  symbolic x coords={Gemma-4,Llama-3-8B,Qwen-2.5-7B},
  xtick=data, tick label style={font=\scriptsize},
  nodes near coords, every node near coord/.append style={font=\scriptsize},
  legend style={font=\scriptsize,at={(0.5,-0.22)},anchor=north,legend columns=2},
  ymajorgrids,
]
\addplot[fill=kcred!60,draw=kcred] coordinates {(Gemma-4,7.40)(Llama-3-8B,11.80)(Qwen-2.5-7B,13.60)};
\addplot[fill=kcgreen!60,draw=kcgreen] coordinates {(Gemma-4,0)(Llama-3-8B,0)(Qwen-2.5-7B,0)};
\legend{Raw RAG hazard,Verified hazard}
\end{axis}
\end{tikzpicture}
\end{center}

Coverage differs across models, but the tested safety boundary remains unchanged.

\begin{pullquote}
{\bfseries\color{kcgreen}The application does not need one perfect language model to become operationally trustworthy. It needs a model-independent boundary around the claims the system is willing to certify.}
\end{pullquote}

\srcfoot{Source: multi-generator robustness evaluation \MEASURED{} \cite{cea_manuscript}.}

\newpage
% ==================== PAGE 21 — BENGALI ROBUSTNESS ====================
\section*{\color{kcgreen}Local language variation should not force unsafe behavior}
\takeaway{When language becomes hard for retrieval, the system uses alternative evidence paths and safe fallback.}

\begin{center}
\begin{tikzpicture}
\begin{axis}[
  width=0.85\linewidth,height=54mm,
  ybar, bar width=16pt, ymin=0, ymax=100,
  ylabel={\footnotesize Coverage (\%)},
  symbolic x coords={Regional dialect,Banglish},
  xtick=data, tick label style={font=\scriptsize},
  nodes near coords, every node near coord/.append style={font=\scriptsize},
  legend style={font=\scriptsize,at={(0.5,-0.22)},anchor=north,legend columns=2},
  ymajorgrids,
]
\addplot[fill=kcearth!55,draw=kcearth] coordinates {(Regional dialect,42.9)(Banglish,41.6)};
\addplot[fill=kcgreen!60,draw=kcgreen] coordinates {(Regional dialect,89.8)(Banglish,90.0)};
\legend{Text-first coverage,Detection-gated coverage}
\end{axis}
\end{tikzpicture}
\end{center}

Detection-gated routing lifts coverage from \textbf{42.9\% to 89.8\%} (regional
dialect) and \textbf{41.6\% to 90.0\%} (Banglish), with zero observed dangerous
acceptance across 4,000 queries.

\vspace{3pt}
The goal is not to claim that Bengali language variation has been solved. The
practical goal is narrower and more useful: when language becomes difficult for
retrieval, the system should have alternative evidence paths and safe fallback
behavior rather than converting linguistic uncertainty into unsupported
agricultural advice.

\srcfoot{Source: dialect-hazard evaluation (E21), 4,000 queries \MEASURED{}.}

\newpage
% ==================== PAGE 22 — EXPERT VALIDATION ====================
\section*{\color{kcgreen}Reviewed by agricultural experts, not only by software tests}
\takeaway{Structured safety behavior was judged by agronomists, not only by the system itself.}

\begin{center}
\begin{tabularx}{\linewidth}{*{4}{>{\centering\arraybackslash}X}}
\bignum{200}{responses reviewed} & \bignum{3}{certified agronomists} & \bignum{4.82/5}{mean correctness} & \bignum{100\%}{chemical safety pass}\\[10mm]
\bignum{98.5\%}{evidence traceability} & \bignum{96.5\%}{deployment approval} & \bignum{0.862}{Gwet's AC1 (safety)} & {}\\
\end{tabularx}
\end{center}

\vspace{4pt}
Automated testing is necessary but not sufficient for agricultural deployment. We
therefore included a double-blind expert review of 200 representative advisory
outputs by three certified agronomists. KrishokChat received a mean correctness
score of 4.82/5, a 100\% chemical-safety pass rate in the evaluated sample, 98.5\%
evidence traceability, and 96.5\% deployment approval, with safety agreement of
Gwet's AC1 = 0.862. These results do not replace field validation, but they
provide a second form of evidence: the structured safety behavior was evaluated by
agricultural specialists rather than judged only by the system itself.

\srcfoot{Source: double-blind expert evaluation \MEASURED{}.}

\newpage
% ==================== PAGE 23 — MULTIMODAL UNCERTAINTY ====================
\section*{\color{kcgreen}When vision is uncertain, uncertainty is passed forward}
\takeaway{A prediction is an input to the advisory system, not permission to act.}

\begin{center}
\begin{tikzpicture}[node distance=4.5mm]
  \node[inbox,text width=34mm] (i) {Image};
  \node[aibox,below=of i,text width=34mm] (c) {Crop prediction};
  \node[aibox,below=of c,text width=34mm] (d) {Disease prediction};
  \node[ambox,below=of d,text width=34mm] (cf) {Confidence};
  \node[ambox,below=of cf,text width=42mm] (q) {Conflict / uncertainty?};
  \draw[fl](i)--(c);\draw[fl](c)--(d);\draw[fl](d)--(cf);\draw[fl](cf)--(q);
  \node[evbox,below left=8mm and -4mm of q,text width=30mm] (no) {NO \textrightarrow{} evidence};
  \node[redbox,below right=8mm and -4mm of q,text width=34mm] (yes) {YES \textrightarrow{} clarification};
  \draw[fl](q)--(no);\draw[fl](q)--(yes);
\end{tikzpicture}
\end{center}

The cross-modal evaluation reports \textbf{100\%} clarification triggering with
\textbf{0\%} unsafe acceptance and no chemical-cocktail delivery across the 100
tested cases.

\begin{pullquote}
{\bfseries\color{kcgreen}A prediction is an input to the advisory system, not permission to act.}
\end{pullquote}

\srcfoot{Source: cross-modal uncertainty evaluation (E31) \MEASURED{}.}

\newpage
% ==================== PAGE 24 — OFFLINE / LAST-MILE ====================
\section*{\color{kcgreen}A correct recommendation still fails if it cannot reach the farmer}
\takeaway{Connectivity is part of the advisory system, not an external assumption.}

\begin{minipage}[t]{0.46\linewidth}
\begin{tikzpicture}[node distance=4mm]
  \node[inbox,text width=38mm] (on) {ONLINE \textrightarrow{} remote advisory};
  \node[ambox,below=of on,text width=38mm] (lim) {LIMITED \textrightarrow{} cached / local evidence};
  \node[evbox,below=of lim,text width=38mm] (off) {OFFLINE \textrightarrow{} verified local pathways};
  \node[greybox,below=of off,text width=38mm] (con) {CONSTRAINED \textrightarrow{} deterministic SMS};
  \draw[fl](on)--(lim);\draw[fl](lim)--(off);\draw[fl](off)--(con);
\end{tikzpicture}
\end{minipage}\hfill
\begin{minipage}[t]{0.5\linewidth}
\begin{tikzpicture}
\begin{axis}[
  width=\linewidth,height=48mm,
  ybar, bar width=13pt, ymin=0, ymax=100,
  ylabel={\footnotesize Delivery (\%)},
  symbolic x coords={15\% loss,30\% loss},
  xtick=data, tick label style={font=\scriptsize},
  nodes near coords, every node near coord/.append style={font=\scriptsize},
  legend style={font=\scriptsize,at={(0.5,-0.24)},anchor=north,legend columns=1},
  ymajorgrids,
]
\addplot[fill=kcred!55,draw=kcred] coordinates {(15\% loss,82.0)(30\% loss,12.8)};
\addplot[fill=kcgreen!60,draw=kcgreen] coordinates {(15\% loss,91.4)(30\% loss,58.1)};
\legend{Cloud-only,Offline-first cache}
\end{axis}
\end{tikzpicture}
\end{minipage}

\vspace{3pt}
\begin{pullquote}
{\bfseries\color{kcgreen}The design treats connectivity as part of the advisory system, not as an external assumption.}
\end{pullquote}

\srcfoot{Controlled network simulation \SIMULATED{}; carrier-level performance remains a deployment validation task.}

\newpage
% ==================== PAGE 25 — CONSTRAINED CHANNELS ====================
\section*{\color{kcgreen}The same safety contract survives the last mile}
\takeaway{We do not ask the language model to decide which safety fields can disappear.}

\begin{center}
\begin{tcolorbox}[colback=kcpanel,colframe=kcgreen,width=0.8\linewidth,arc=2pt]
\footnotesize\ttfamily
[KrishokChat] Alu / Late blight\\
AI: Mancozeb 80\% WP\\
Dose: 2 g/L water\\
Interval: 7-10 d\ \ PHI: 7 d\\
PPE: mask+gloves\ \ Call 16123
\end{tcolorbox}
\end{center}

\begin{center}\footnotesize\emph{Deterministic SMS template: 102--115 characters, critical dosage/interval/PHI fields preserved.}\end{center}

\vspace{3pt}
\begin{center}
\begin{tabular}{lrr}
\toprule
\textbf{Channel composition} & \textbf{Critical-hazard} & \textbf{Injection leakage}\\
\midrule
LLM-composed control & 64.4\% & 509 / 1{,}400\\
Deterministic template & --- & 0 / 1{,}400\\
\bottomrule
\end{tabular}
\end{center}

\begin{pullquote}
{\bfseries\color{kcgreen}When a message must be compressed, we do not ask the language model to decide which safety fields can disappear. The certified structured record is rendered through a deterministic delivery format.}
\end{pullquote}

\srcfoot{Source: GSM composition (E15) and SMS injection (E22) \MEASURED{}.}

\newpage
% ==================== PAGE 26 — EFFICIENCY ====================
\section*{\color{kcgreen}Use generation where it adds value; do not pay for it when the answer is known}
\takeaway{Stop using a generative model as the default authority for every agricultural question.}

\begin{center}
\begin{tikzpicture}
\begin{axis}[
  width=0.92\linewidth,height=58mm,
  xbar, xmin=0, xmax=45,
  bar width=8pt,
  enlarge y limits=0.18,
  xlabel={\footnotesize Share of 5,000-query workload (\%)},
  symbolic y coords={T4 Refusal,T0 Safety,T2 Templated facts,T1 Detect$\to$Facts,T3 LLM},
  ytick=data,
  nodes near coords,
  every node near coord/.append style={font=\scriptsize},
  tick label style={font=\scriptsize},
  xmajorgrids,
]
\addplot[fill=kcgreen!55,draw=kcgreen] coordinates {
  (4.08,T4 Refusal)(6.40,T0 Safety)(18.82,T2 Templated facts)(32.22,T1 Detect$\to$Facts)(38.48,T3 LLM)};
\end{axis}
\end{tikzpicture}
\end{center}

\begin{center}
\begin{tabularx}{0.8\linewidth}{*{2}{>{\centering\arraybackslash}X}}
\bignum{51.04\%}{workload on deterministic advisory paths} &
\bignum{61.52\%}{reaches a zero-LLM terminal path (incl. safety/refusal)}\\
\end{tabularx}
\end{center}

\vspace{2pt}
Weighted latency drops from \textbf{1,247.93\,ms} (LLM-default) to
\textbf{546.15\,ms} (tiered routing).

\begin{pullquote}
{\bfseries\color{kcgreen}The point is not to eliminate LLMs. The point is to stop using a generative model as the default authority for every agricultural question.}
\end{pullquote}

\srcfoot{Source: 5,000-query workload distribution \MEASURED{} \cite{cea_manuscript}.}

\newpage
% ==================== PAGE 27 — PRODUCTION ====================
\section*{\color{kcgreen}The software architecture was designed to evolve into a service}
\takeaway{New models, evidence packs, crops and channels can be added without a rebuild.}

\begin{minipage}[t]{0.4\linewidth}
\begin{tcolorbox}[colback=black!4,colframe=kcgrey,arc=2pt]
\footnotesize\ttfamily
frontend/\\ backend/\\ capstone/\\ dataset\_release/\\ demo-assets/\\ deploy/\\ docs/\\ paper/\\ production/future\_plan/\\ research\_artifacts/\\ scripts/\\ tools/
\end{tcolorbox}
\end{minipage}\hfill
\begin{minipage}[t]{0.56\linewidth}
KrishokChat is not tied to one model provider, one retrieval implementation, or
one frontend workflow. The codebase separates application logic from
infrastructure adapters, keeps research artifacts separate from runtime assets,
and provides explicit deployment, audit and testing layers.

\vspace{4pt}
This matters commercially because the agricultural knowledge, models and external
services will change over time. The platform therefore accommodates new models,
new evidence packs, new crops and new delivery channels without forcing a
complete rebuild.
\end{minipage}

\vspace{4pt}
\begin{center}
\begin{tabularx}{\linewidth}{*{4}{>{\centering\arraybackslash}X}}
\bignum{558}{backend tests passed} & \bignum{8}{skipped} & \bignum{0}{failed} & \bignum{50/50}{golden replay}\\
\end{tabularx}
\end{center}

\srcfoot{Source: repository (193 commits) and CI baseline \MEASURED{} \cite{krishokchat_repo}. Frontend build green.}

\newpage
% ==================== PAGE 28 — KNOWLEDGE GROWTH ====================
\section*{\color{kcgreen}The long-term asset is the governed agricultural knowledge layer}
\takeaway{Knowledge improvement is an update to the platform, not a rebuild of the application.}

\begin{center}
\begin{tikzpicture}[node distance=6mm]
  \node[evbox,text width=32mm] (s) {Authoritative source};
  \node[greybox,right=of s,text width=28mm] (e) {Fact extraction};
  \node[ambox,right=of e,text width=26mm] (v) {Validation};
  \node[evbox,below=6mm of v,text width=32mm] (p) {Versioned knowledge pack};
  \node[greybox,left=of p,text width=30mm] (a) {Advisory system};
  \draw[fl](s)--(e);\draw[fl](e)--(v);\draw[fl](v)--(p);\draw[fl](p)--(a);
\end{tikzpicture}
\end{center}

\vspace{3pt}
A targeted knowledge intervention in one chili-anthracnose cluster increased
coverage from \textbf{58.7\% to 64.2\%} in the evaluated sample --- two
BARI-verified facts, added in 18 minutes, over a 55-query cluster.

\begin{pullquote}
{\bfseries\color{kcgreen}Knowledge improvement becomes an update to the platform, not a rebuild of the application.}
\end{pullquote}

\srcfoot{Source: knowledge-growth evaluation (E24), single cluster \MEASURED{}.}

\newpage
% ==================== PAGE 29 — PROVENANCE ====================
\section*{\color{kcgreen}Every important agricultural claim should remain traceable}
\takeaway{The objective is to preserve the path from source to the certified fact.}

\begin{center}
\begin{tikzpicture}[node distance=4mm]
  \node[evbox,text width=40mm] (a) {Source publication};
  \node[evbox,below=of a,text width=40mm] (b) {Semantic unit};
  \node[evbox,below=of b,text width=40mm] (c) {Structured fact};
  \node[greybox,below=of c,text width=40mm] (d) {Advisory claim};
  \node[ambox,below=of d,text width=40mm] (e) {Verification};
  \node[evbox,below=of e,text width=40mm] (f) {Delivered response};
  \draw[fl](a)--(b);\draw[fl](b)--(c);\draw[fl](c)--(d);\draw[fl](d)--(e);\draw[fl](e)--(f);
\end{tikzpicture}
\end{center}

\vspace{3pt}
In a 1,000-mutation provenance test, \textbf{1,000/1,000} tampering attempts were
detected. The hash-chained evidence pack carries only a \textbf{$\sim$791-byte}
delta against a 10,985-byte full pack, with \textbf{0.1919\,ms} p95 verification
in the tested client workflow.

\begin{pullquote}
{\bfseries\color{kcgreen}The objective is not merely to cite a source. It is to preserve the path from the source to the specific agricultural fact that the system was willing to certify.}
\end{pullquote}

\srcfoot{Source: provenance/cache evaluation (E23) \MEASURED{}.}

\newpage
% ==================== PAGE 30 — RESILIENCE MATRIX ====================
\section*{\color{kcgreen}When something fails, the system degrades deliberately}
\takeaway{Every failure mode has a defined, safe preferred behavior.}

\begin{center}
\begin{tabularx}{\linewidth}{Xl}
\toprule
\textbf{Failure} & \textbf{Preferred behavior}\\
\midrule
Unsafe query & Block\\
Missing context & Clarify\\
Unsupported fact & Abstain\\
Evidence conflict & Reject / escalate\\
Stale source & Invalidate\\
Vision / text conflict & Clarify\\
Remote LLM unavailable & Local / deterministic path\\
Dense retrieval unavailable & Alternate retrieval\\
Network unavailable & Cached / local path\\
SMS too long & Deterministic refusal / recomposition\\
\bottomrule
\end{tabularx}
\end{center}

\vspace{4pt}
Each row is a designed operating state rather than an unhandled exception. This is
what makes the project operational, not merely intelligent: the failure surfaces
were enumerated, and each has a defined safe response.

\srcfoot{Source: KrishokChat operating-state design \cite{cea_manuscript,krishokchat_repo}.}

\newpage
% ==================== PAGE 31 — VALIDATION MATRIX ====================
\section*{\color{kcgreen}What has been tested}
\takeaway{The evaluation program is broader than any single benchmark because the system has multiple failure surfaces.}

\begin{minipage}[t]{0.48\linewidth}
\textbf{\color{kcgreen}Safety}\par\footnotesize
relational misbinding \textbullet{} slot ablation \textbullet{} counterfactual
evidence \textbullet{} prompt injection \textbullet{} retrieval poisoning
\textbullet{} stale evidence \textbullet{} source fragmentation \textbullet{}
model variation

\vspace{5pt}
\textbf{\color{kcgreen}Language}\par\footnotesize
standard Bengali \textbullet{} farmer language \textbullet{} regional dialect
\textbullet{} Banglish \textbullet{} normalization

\vspace{5pt}
\textbf{\color{kcgreen}Multimodal}\par\footnotesize
crop classification \textbullet{} disease classification \textbullet{} cross-modal
conflict \textbullet{} confidence uncertainty
\end{minipage}\hfill
\begin{minipage}[t]{0.48\linewidth}
\textbf{\color{kcgreen}Deployment}\par\footnotesize
latency \textbullet{} local inference \textbullet{} cache \textbullet{} degraded
networks \textbullet{} GSM composition \textbullet{} provenance updates
\textbullet{} knowledge maintenance

\vspace{5pt}
\textbf{\color{kcgreen}Human}\par\footnotesize
expert correctness \textbullet{} safety \textbullet{} evidence traceability
\textbullet{} deployment approval
\end{minipage}

\vspace{6pt}
\begin{pullquote}
{\bfseries\color{kcgreen}The evaluation program is broader than any single benchmark because the system has multiple failure surfaces.}
\end{pullquote}

\newpage
% ==================== PAGE 32 — FIVE CONCLUSIONS ====================
\section*{\color{kcgreen}Five conclusions we can support today}
\takeaway{Each conclusion is tied to a specific body of evidence.}

\newcommand{\concl}[3]{\begin{tcolorbox}[colback=kcpanel,colframe=kcgreen,arc=2pt,left=8pt,right=8pt,top=4pt,bottom=4pt]
{\bfseries\color{kcgreen}#1}\quad #2\par\vspace{2pt}{\footnotesize\color{kcgrey}#3}\end{tcolorbox}\vspace{3pt}}

\concl{01}{Agricultural generation cannot be treated as a purely linguistic problem.}{Supported by KrishokChat v2.}
\concl{02}{Retrieval quality changes with farmer language and query condition.}{Supported by the retrieval study.}
\concl{03}{Evidence binding must respect agricultural relationships, not only lexical overlap.}{Supported by counterfactual and slot-ablation results.}
\concl{04}{The system can trade answer coverage for safer behavior instead of guessing.}{Supported by retrieval-degradation and selective behavior.}
\concl{05}{The application has a viable engineering path from prototype toward deployment.}{Supported by codebase, deployment layers and test infrastructure --- not by field adoption.}

\begin{center}{\large\bfseries\color{kcgreen}These findings led us to build KrishokChat as a platform, not a chatbot.}\end{center}

\newpage
% ==================== PAGE 33 — COMMERCIAL MODEL ====================
\section*{\color{kcgreen}Designed as shared advisory infrastructure}
\takeaway{The farmer is the beneficiary; institutions can finance deployment.}

\begin{center}
\begin{tikzpicture}[node distance=8mm and 12mm]
  \node[evbox,text width=40mm] (kc) {\bfseries KRISHOKCHAT};
  \node[inbox,below left=10mm and 8mm of kc,text width=26mm] (farm) {Farmers};
  \node[greybox,below=10mm of kc,text width=26mm] (ngo) {NGOs};
  \node[greybox,below right=10mm and 8mm of kc,text width=26mm] (gov) {Government};
  \draw[fl](kc)--(farm);\draw[fl](kc)--(ngo);\draw[fl](kc)--(gov);
  \node[greybox,below=8mm of ngo,text width=30mm] (agri) {Agribusiness};
  \node[greybox,below=8mm of agri,text width=34mm] (ext) {Extension network};
  \draw[fl](ngo)--(agri);\draw[fl](agri)--(ext);
\end{tikzpicture}
\end{center}

\begin{minipage}[t]{0.32\linewidth}
\begin{card}{Farmer-facing layer}\footnotesize Free / basic access.\end{card}
\end{minipage}\hfill
\begin{minipage}[t]{0.32\linewidth}
\begin{card}{Institutional layer}\footnotesize Dashboards, APIs, knowledge packs, deployment support.\end{card}
\end{minipage}\hfill
\begin{minipage}[t]{0.32\linewidth}
\begin{card}{Public-service layer}\footnotesize Extension integration, escalation, broadcast advisory.\end{card}
\end{minipage}

\vspace{5pt}
The most realistic commercialization path is not to assume that individual
smallholder farmers will finance the entire service directly. KrishokChat is
better positioned as shared advisory infrastructure: the farmer is the
beneficiary, while institutions can finance deployment, support, knowledge
integration and operational services. The same core platform can support multiple
channels without creating separate products for each organization.

\srcfoot{Production planning identifies institutional / B2B / B2G deployment as the more durable path \cite{krishokchat_repo}.}

\newpage
% ==================== PAGE 34 — MARKET NEED ====================
\section*{\color{kcgreen}The demand already exists; the missing layer is intelligent delivery}
\takeaway{KrishokChat adds a digital intelligence layer to an existing ecosystem; it does not replace extension institutions.}

Bangladesh does not need to invent demand for agricultural advice. Farmers already
seek support through extension officers, agricultural call centres (the national
helpline \textbf{16123}), weather and advisory services (BAMIS), local experts and
informal information networks.

\vspace{4pt}
The opportunity for KrishokChat is to add a digital intelligence layer to that
existing ecosystem: understanding farmer-language questions, combining text and
image evidence, grounding responses in agricultural knowledge, enforcing safety
checks, and routing unresolved cases toward human or institutional support. This
is much stronger than pretending KrishokChat replaces existing extension
institutions.

\vspace{4pt}
\begin{center}
\begin{tabularx}{0.85\linewidth}{*{3}{>{\centering\arraybackslash}X}}
\begin{card}{BAMIS}\footnotesize Agricultural advisory information service.\end{card} &
\begin{card}{16123}\footnotesize National agricultural call center.\end{card} &
\begin{card}{Extension}\footnotesize Field officer network (SAAOs).\end{card}\\
\end{tabularx}
\end{center}

\srcfoot{Sources: existing Bangladesh advisory infrastructure \cite{innovation_fair}.}

\newpage
% ==================== PAGE 35 — SOCIAL INCLUSION ====================
\section*{\color{kcgreen}Access should not depend on English, perfect spelling or reliable broadband}
\takeaway{Inclusion is implemented as product behavior rather than a statement of intent.}

\begin{center}
\begin{tabularx}{\linewidth}{*{5}{>{\centering\arraybackslash}X}}
\begin{card}{Bengali-first}\footnotesize Plain Bengali interaction.\end{card} &
\begin{card}{Dialect-aware}\footnotesize Multiple regional registers.\end{card} &
\begin{card}{Banglish-tolerant}\footnotesize Romanized input handled.\end{card} &
\begin{card}{Low-connectivity}\footnotesize Cache and offline pathways.\end{card} &
\begin{card}{Voice-ready}\footnotesize Designed for future voice input.\end{card}\\
\end{tabularx}
\end{center}

\vspace{4pt}
\begin{pullquote}
{\bfseries\color{kcgreen}Inclusion is implemented as product behavior rather than a statement of intent.}
\end{pullquote}

Concretely: large touch targets, plain Bengali phrasing, image-based input,
progressive clarification, offline / cache pathways, and a future voice-input
affordance. The frontend roadmap considers WCAG-oriented touch targets, Bengali
font optimization, PWA / offline UX, on-device diagnosis and voice.

\srcfoot{Source: frontend roadmap \cite{krishokchat_repo}.}

\newpage
% ==================== PAGE 36 — ENVIRONMENT ====================
\section*{\color{kcgreen}Better information can also mean more responsible intervention}
\takeaway{Treatment recommendations are made conditional on evidence, supporting IPM where authoritative guidance exists.}

We do not claim that KrishokChat reduces pesticide usage nationally. Instead, the
system is designed to make treatment recommendations conditional on evidence and
to support integrated pest management (IPM) and non-chemical alternatives where
authoritative guidance exists.

\vspace{5pt}
\begin{center}
\begin{tabularx}{\linewidth}{*{5}{>{\centering\arraybackslash}X}}
\begin{card}{Dosage verification}\footnotesize Bounded, evidence-checked doses.\end{card} &
\begin{card}{PHI}\footnotesize Pre-harvest interval enforced.\end{card} &
\begin{card}{Regulatory state}\footnotesize Approved / restricted checked.\end{card} &
\begin{card}{IPM}\footnotesize Integrated management surfaced.\end{card} &
\begin{card}{Alternatives}\footnotesize Non-chemical options where supported.\end{card}\\
\end{tabularx}
\end{center}

\vspace{4pt}
This is where ethics connects to actual architecture: the same typed treatment
contract that enforces safety also makes it possible to prefer lower-risk,
evidence-supported interventions.

\newpage
% ==================== PAGE 37 — SDGs ====================
\section*{\color{kcgreen}The innovation contributes directly to four development priorities}
\takeaway{Focused SDG alignment, not a scattered checklist.}

\begin{center}
\begin{tabularx}{\linewidth}{*{4}{>{\centering\arraybackslash}X}}
\begin{card}{SDG 2 --- Zero Hunger}\footnotesize Agricultural knowledge, extension, resilience.\end{card} &
\begin{card}{SDG 9 --- Industry \& Infrastructure}\footnotesize Local digital infrastructure.\end{card} &
\begin{card}{SDG 12 --- Responsible Production}\footnotesize Evidence-bound intervention, safer chemical guidance.\end{card} &
\begin{card}{SDG 13 --- Climate Action}\footnotesize Weather-aware / proactive advisory pathways.\end{card}\\
\end{tabularx}
\end{center}

\vspace{5pt}
SDG 2 explicitly targets smallholder productivity and income, sustainable
agriculture, resilience, extension and technological development. Bangladesh's own
SDG commitment connects digital agriculture to productivity, food security and the
empowerment of rural farmers. KrishokChat aligns with these priorities by being a
locally grounded, safety-aware agricultural intelligence layer.

\srcfoot{Sources: UN SDG 2 targets \cite{sdg2}; Bangladesh SDG commitments \cite{innovation_fair}.}

\newpage
% ==================== PAGE 38 — LEADERSHIP / TIMELINE ====================
\section*{\color{kcgreen}Developed as a continuous research-to-engineering cycle}
\takeaway{Each stage produced an artifact the next stage depended on.}

\begin{center}
\begin{tikzpicture}[node distance=3mm]
  \node[inbox,text width=30mm] (a) {Field visit};
  \node[inbox,below=of a,text width=30mm] (b) {Problem identification};
  \node[evbox,below=of b,text width=30mm] (c) {Dataset / benchmark};
  \node[evbox,below=of c,text width=30mm] (d) {Retrieval analysis};
  \node[greybox,below=of d,text width=30mm] (e) {System development};
  \node[ambox,below=of e,text width=30mm] (f) {Safety research};
  \node[greybox,below=of f,text width=30mm] (g) {Production hardening};
  \node[evbox,below=of g,text width=30mm] (h) {Field deployment target};
  \draw[fl](a)--(b);\draw[fl](b)--(c);\draw[fl](c)--(d);\draw[fl](d)--(e);
  \draw[fl](e)--(f);\draw[fl](f)--(g);\draw[fl](g)--(h);
\end{tikzpicture}
\end{center}

\begin{center}
\begin{tabularx}{\linewidth}{*{4}{>{\centering\arraybackslash}X}}
\begin{card}{Research}\footnotesize Benchmark, retrieval, safety studies.\end{card} &
\begin{card}{ML}\footnotesize Models, vision, verification.\end{card} &
\begin{card}{Backend}\footnotesize Pipeline, safety, audit.\end{card} &
\begin{card}{Frontend}\footnotesize PWA, Bengali UX.\end{card}\\[4pt]
\begin{card}{Infrastructure}\footnotesize Deployment, testing.\end{card} &
\begin{card}{Field coordination}\footnotesize Farmer interviews.\end{card} &
\begin{card}{Validation}\footnotesize Expert review.\end{card} & {}\\
\end{tabularx}
\end{center}

\newpage
% ==================== PAGE 39 — DIFFERENTIATION ====================
\section*{\color{kcgreen}What separates KrishokChat from a feature demo}
\takeaway{Problem-first, evidence-bound, and operationally engineered.}

\begin{center}
\begin{tabularx}{\linewidth}{XX}
\toprule
\textbf{Typical prototype} & \textbf{KrishokChat}\\
\midrule
Model-first & Problem-first\\
One model & Multiple controlled capabilities\\
Answer-only & Answer + evidence + verification\\
Online assumed & Degraded connectivity considered\\
Static data & Governed evidence layer\\
No audit & Audit trail\\
No structured failure state & Clarify / refuse / escalate\\
One provider & Replaceable model providers\\
Demo code & Modular deployment-oriented architecture\\
Evaluation = accuracy & Safety + coverage + latency + reliability\\
\bottomrule
\end{tabularx}
\end{center}

\newpage
% ==================== PAGE 40 — ROADMAP ====================
\section*{\color{kcgreen}The next step is not another prototype --- it is a field deployment cycle}
\takeaway{From validated system to national service layer, in defined phases.}

\begin{center}
\begin{tikzpicture}[node distance=5mm]
  \node[evbox,text width=34mm] (c) {CURRENT: Validated system};
  \node[greybox,below=of c,text width=34mm] (p1) {PHASE 1: Controlled pilot};
  \node[greybox,below=of p1,text width=34mm] (p2) {PHASE 2: Farmer + extension evaluation};
  \node[greybox,below=of p2,text width=34mm] (p3) {PHASE 3: Institutional deployment};
  \node[greybox,below=of p3,text width=34mm] (p4) {PHASE 4: Regional scale};
  \node[evbox,below=of p4,text width=34mm] (p5) {PHASE 5: National service layer};
  \draw[fl](c)--(p1);\draw[fl](p1)--(p2);\draw[fl](p2)--(p3);\draw[fl](p3)--(p4);\draw[fl](p4)--(p5);
\end{tikzpicture}
\end{center}

\begin{tabularx}{\linewidth}{lX}
\textbf{Phase 1} & Deploy to a controlled farmer cohort.\\
\textbf{Phase 2} & Measure expert agreement, useful coverage, safety, referral and user behavior.\\
\textbf{Phase 3} & Integrate with institutional workflows.\\
\textbf{Phase 4} & Expand crop / knowledge / service coverage.\\
\textbf{Phase 5} & Scale through public and private agricultural channels.\\
\end{tabularx}

\newpage
% ==================== PAGE 41 — WHAT FUNDING UNLOCKS ====================
\section*{\color{kcgreen}What support changes}
\takeaway{The next investment reduces the gap between controlled validation and field operation.}

\begin{center}
\begin{tabularx}{\linewidth}{*{2}{>{\raggedright\arraybackslash}X}}
\begin{card}{Field validation}\footnotesize Farmer pilots, extension evaluation, outcome measurement.\end{card} &
\begin{card}{Knowledge scale}\footnotesize Expert-reviewed agricultural knowledge packs, regulatory updates.\end{card}\\[4pt]
\begin{card}{Deployment}\footnotesize Hosting, edge / offline, monitoring, service reliability.\end{card} &
\begin{card}{Ecosystem}\footnotesize Institutional integration, extension workflow, partnership development.\end{card}\\
\end{tabularx}
\end{center}

\vspace{5pt}
The core system does not need to be invented from scratch. The next investment is
in reducing the gap between controlled validation and field operation. Funding
would support expert-reviewed knowledge expansion, farmer and extension pilots,
production infrastructure, multimodal and offline deployment, operational
monitoring, and institutional integration. The objective is to measure real-world
value under real conditions and build the evidence required for responsible scale.

\newpage
% ==================== PAGE 42 — SUCCESS METRICS ====================
\section*{\color{kcgreen}We will measure deployment by outcomes, not downloads}
\takeaway{Success is defined by safe, correct, resolved advisory cases.}

\begin{center}
\begin{tabularx}{\linewidth}{*{5}{>{\centering\arraybackslash}X}}
\begin{card}{Safe advisory rate}\footnotesize\end{card} &
\begin{card}{Correct advisory rate}\footnotesize\end{card} &
\begin{card}{Appropriate abstention}\footnotesize\end{card} &
\begin{card}{Farmer task completion}\footnotesize\end{card} &
\begin{card}{Extension referral resolution}\footnotesize\end{card}\\
\end{tabularx}
\end{center}

\vspace{5pt}
Longer-term outcome measures include crop-loss / yield outcomes, chemical-use
outcomes, user trust, retention, and cost per successfully resolved case. These
are field-validation targets: they define what deployment must prove, and they are
deliberately stated as outcomes rather than vanity metrics.

\newpage
% ==================== PAGE 43 — EVIDENCE STATUS ====================
\section*{\color{kcgreen}What we know now, and what deployment must still prove}
\takeaway{We know exactly where the evidence currently ends.}

\begin{minipage}[t]{0.48\linewidth}
\textbf{\color{kcgreen}Established today}\par\footnotesize
benchmark resources \textbullet{} retrieval findings \textbullet{} system
architecture \textbullet{} safety tests \textbullet{} expert review \textbullet{}
software verification \textbullet{} network simulation \textbullet{} local model
path \textbullet{} production structure
\end{minipage}\hfill
\begin{minipage}[t]{0.48\linewidth}
\textbf{\color{kcearth}Next validation}\par\footnotesize
field outcomes \textbullet{} carrier-level delivery \textbullet{} farmer adoption
\textbullet{} extension workload \textbullet{} real referral follow-through
\textbullet{} long-term trust \textbullet{} yield / chemical-use impact
\end{minipage}

\vspace{8pt}
\begin{center}\footnotesize
\begin{tabular}{lll}
\toprule
\textbf{Evidence} & \textbf{Status} & \textbf{Meaning}\\
\midrule
v2 benchmark & Established & arXiv resource\\
Retrieval study & Established & arXiv study\\
Expert review & Measured & 200 outputs\\
Safety attack suite & Measured & specified benchmark\\
Counterfactual verification & Measured & 2,000 pairs\\
Network resilience & Simulated & emulator\\
GSM composition & Tested & software composition\\
Carrier delivery & Future & not measured\\
Farmer outcome & Future & field study needed\\
National impact & Future & not yet demonstrated\\
\bottomrule
\end{tabular}
\end{center}

\begin{pullquote}
{\bfseries\color{kcgreen}We know exactly where the evidence currently ends.}
\end{pullquote}

\newpage
% ==================== PAGE 44 — RESEARCH PORTFOLIO ====================
\section*{\color{kcgreen}Backed by a continuous research programme}
\takeaway{Four connected research contributions underpin the platform.}

\begin{tabularx}{\linewidth}{>{\raggedright\arraybackslash}X}
\begin{card}{Research 01 --- Provenance-Traceable Bengali Agricultural Benchmark}\footnotesize
\textbf{Contribution:} knowledge / evaluation foundation.\quad\textbf{Status:} arXiv preprint (current) \cite{krishokchat_v2}.\end{card}\\[4pt]
\begin{card}{Research 02 --- Where Does Retrieval Fail?}\footnotesize
\textbf{Contribution:} retrieval failure analysis.\quad\textbf{Status:} arXiv:2608.14886 \cite{retrieval_study}.\end{card}\\[4pt]
\begin{card}{Research 03 --- Bounded-Authority Agricultural Advisory (CEA manuscript)}\footnotesize
\textbf{Contribution:} system reliability, routing, verification, deployment architecture.\quad\textbf{Status:} in preparation \cite{cea_manuscript}.\end{card}\\[4pt]
\begin{card}{Research 04 --- KrishokChat System Demonstration}\footnotesize
\textbf{Contribution:} integrated public system.\quad\textbf{Status:} demonstration / submission.\end{card}
\end{tabularx}

\newpage
% ==================== PAGE 45 — LINKS ====================
\section*{\color{kcgreen}Explore the work}
\takeaway{The platform and its research are openly accessible.}

\begin{center}
\begin{tabularx}{\linewidth}{*{4}{>{\centering\arraybackslash}X}}
\figslot{figures/qr_live.png}{30mm}{} &
\figslot{figures/qr_github.png}{30mm}{} &
\figslot{figures/qr_v2.png}{30mm}{} &
\figslot{figures/qr_retrieval.png}{30mm}{}\\
{\bfseries Live system}\par{\footnotesize Try KrishokChat} &
{\bfseries Source code}\par{\footnotesize GitHub repository} &
{\bfseries Research}\par{\footnotesize KrishokChat v2} &
{\bfseries Retrieval study}\par{\footnotesize Where Does Retrieval Fail?}\\
\end{tabularx}
\end{center}

\vspace{6pt}
\begin{center}\small\textbf{Primary repository:} \url{https://github.com/RaiyaanReza/KrishokChat-Agricultural-Advisory-System}\end{center}

\srcfoot{The repository is public and exposes backend, frontend, deployment, research artifacts, production plans and experiment materials \cite{krishokchat_repo}. QR placeholders swap to real codes when generated.}

\newpage
% ==================== PAGE 46 — CLOSING ====================
\thispagestyle{empty}
\vspace{6mm}
\begin{center}{\fontsize{26}{30}\selectfont\bfseries\color{kcgreen}From field problem to trusted agricultural service}\end{center}

\vspace{4mm}
\figslot{figures/closing_system.png}{60mm}{}

\vspace{5mm}
KrishokChat began with a practical question from the field:

\vspace{2mm}
\begin{center}\itshape\large How can agricultural technology become genuinely useful to a farmer when
language, evidence, uncertainty and connectivity all matter at the same time?\end{center}

\vspace{2mm}
We answered that question by building the research foundation, measuring where
current AI fails, redesigning the advisory architecture around those failures, and
engineering the result into a working platform. The next step is to take that
system into the field, measure its real value, and build the institutional pathway
for responsible scale.

\vspace{6mm}
\begin{center}{\Large\bfseries\color{kcgreen}Research \textrightarrow{} Evidence \textrightarrow{} System \textrightarrow{} Field}\end{center}

\vspace{8mm}
\begin{pullquote}
{\bfseries\color{kcgreen}KrishokChat is no longer only a question of whether an AI model can answer an agricultural question. The project is now about whether Bangladesh can build a locally grounded, evidence-aware, safe and deployable intelligence layer for agricultural decision support.}
\end{pullquote}

\vspace{6mm}
\begin{center}{\Large\bfseries\color{kcearth}The next experiment is in the field.}\end{center}

\newpage
% ==================== APPENDICES ====================
\section*{\color{kcgreen}Appendix A --- Experiment catalogue (grouped by question)}
{\footnotesize Internal experiment IDs shown for traceability only.}

\begin{tabularx}{\linewidth}{lX}
\toprule
\textbf{Theme} & \textbf{Internal experiments}\\
\midrule
Knowledge \& evidence & E19 graph traversal, E23 provenance/cache, E24 knowledge growth, E30 temporal validity, E36 source fragmentation\\
Routing \& retrieval & E17 detection gating, E25 compact router, E32 oracle routing, E33 wrong high-confidence routing, E35 normalization, E12 retrieval degradation\\
Safety \& verification & E02 misbinding, E03 slot ablation, E04 calibration, E05 counterfactual, E07/E08 security, E10 failure taxonomy, E11 generator robustness, E12 poisoning, E22 SMS injection\\
Language \& multimodality & E06 dialect, E21 dialect hazard, E31 multimodal uncertainty, E37 clarification\\
Delivery \& deployment & E09 runtime, E14 network, E15 GSM, E18 LLM dependency, E20 cost, E38 escalation, E39 IPM\\
\bottomrule
\end{tabularx}

\subsection*{Appendix B --- Key empirical values}
\begin{longtable}{lrl}
\toprule
\textbf{Metric} & \textbf{Value} & \textbf{Status}\\
\midrule
\endhead
Core benchmark instances & 85{,}979 & Established\\
Real-world farmer queries & 1{,}000 & Established\\
Semantic knowledge units & 2{,}946 & Established\\
Source publications & 284 & Established\\
Hybrid RRF R@10 & 0.539 & Measured\\
Counterfactual certification (typed) & 0.00\% & Measured\\
Counterfactual certification (Vanilla RAG) & 72.65\% & Measured\\
Slot ablation --- dosage bounds & +31.6 pp & Measured\\
Retrieval degradation hazard (KrishokChat) & 0.0\% & Measured\\
Verified hazard (all generators) & 0.00\% & Measured\\
Dialect coverage (gated) & 89.8 / 90.0\% & Measured\\
Expert mean correctness & 4.82 / 5 & Measured\\
Deployment approval & 96.5\% & Measured\\
Gwet's AC1 (safety) & 0.862 & Measured\\
Cross-modal clarification trigger & 100\% & Measured\\
Deterministic advisory workload & 51.04\% & Measured\\
Zero-LLM terminal handling & 61.52\% & Measured\\
Weighted latency (tiered vs LLM) & 546.15 / 1247.93 ms & Measured\\
Offline delivery @30\% loss & 58.1\% vs 12.8\% & Simulated\\
SMS injection leakage (deterministic) & 0 / 1{,}400 & Measured\\
Provenance mutation detection & 1{,}000 / 1{,}000 & Measured\\
Backend tests & 558 pass / 8 skip / 0 fail & Measured\\
\bottomrule
\end{longtable}

\subsection*{Appendix C --- Evidence discipline}
Every quantitative claim in this dossier carries an evidence label:
\MEASURED{} real experimental measurement; \SIMULATED{} software / network
simulation; \ESTABLISHED{} published research resource; \TESTED{} software
composition; \PROJECTED{} scenario or model assumption. Simulated network results
are not carrier-level measurements. Field outcomes, farmer adoption and national
impact are explicitly future-validation items.

\newpage
% ==================== REFERENCES ====================
\section*{\color{kcgreen}References}
\bibliographystyle{plain}
\bibliography{references}

\end{document}

